\section{Saturation of the area law}

\begin{definition}[Block entropy of an MPDO]\label{def:block_entropy}
    \lean{MPOTensor.blockEntropy}
    \leanok
    \uses{def:mpo_tensor, def:von_neumann_entropy}
    For a fixed system size $N$, a block length $L\le N$, and an MPO tensor $M$
    such that $\rho^{(N)}(M)$ is positive semidefinite, let
    \begin{align}
        \sigma^{(N)}(M)
        &=\frac{\rho^{(N)}(M)}{\tr[\rho^{(N)}(M)]}, \notag\\
        \rho_L^{(N)}(M)
        &=\tr_{N\setminus L}\![\sigma^{(N)}(M)]
        \notag
    \end{align}
    denote the normalized state and the reduced state of the first $L$ of $N$
    spins. The \emph{$L$-block entropy} is
    $S_L^{(N)}(M)=S(\rho_L^{(N)}(M))$.
\end{definition}

\begin{definition}[Mutual information of an MPDO]\label{def:mutual_info_chain}
    \lean{MPOTensor.mutualInfoChain, MPOTensor.mutualInfoChain_eq}
    \leanok
    \uses{def:block_entropy}
    For the same fixed system size $N$, block length $L\le N$, and positive
    semidefiniteness of $\rho^{(N)}(M)$, the \emph{mutual information} between
    a block of $L$ spins and the rest is
    $I_L=S_L^{(N)}(M)+S_{N-L}^{(N)}(M)-S_N^{(N)}(M)$.
\end{definition}

\begin{definition}[Saturation of the area law]\label{def:is_sal}
    \lean{MPOTensor.IsSAL}
    \leanok
    \uses{def:mutual_info_chain}
    An MPO tensor $M$ that generates an MPDO and whose positive-length
    finite-chain operators $\rho^{(N)}(M)$ all have nonzero trace (so the normalized states are
    well-defined) \emph{saturates the area law} if
    $I_L=I_{L+1}$ for all $1\le L<\lfloor N/2\rfloor$ and all $N$; that is,
    the mutual information is constant,
    $I_1=I_2=\cdots=I_{\lfloor N/2\rfloor}$.
    This is \cite[Definition~4.6, line~811]{Cirac2016MPDO_arXiv}.
\end{definition}

\begin{lemma}[Constancy on a finite interval]
    \label{lem:finite_interval_constancy}
    \lean{eq_of_forall_succ_eq_of_mem_Icc}
    \leanok
    Let $a,n\in\mathbb N$, and let $f$ be defined on the integer interval
    $[0,n]$. If $f(m)=f(m+1)$ for every $a\le m<n$, then $f(i)=f(j)$ for all
    $i,j\in[a,n]$.
\end{lemma}
\begin{proof}\leanok
    Compose the successive equalities from the smaller index to the larger.
\end{proof}

\begin{lemma}[Saturation telescopes to a constant mutual information]
    \label{lem:sal_const}
    \lean{MPOTensor.mutualInfoChain_eq_of_isSAL}
    \leanok
    \uses{def:is_sal, lem:finite_interval_constancy}
    If an MPO tensor $M$ saturates the area law, then the mutual informations
    coincide, $I_L=I_{L'}$, for all
    $1\le L,L'\le\lfloor N/2\rfloor$.
\end{lemma}
\begin{proof}\leanok
    \uses{def:is_sal, lem:finite_interval_constancy}
    The defining condition supplies $I_m=I_{m+1}$ throughout the interval
    $1\le m<\lfloor N/2\rfloor$. Apply
    Lemma~\ref{lem:finite_interval_constancy}.
\end{proof}

\begin{theorem}[Tripartite marginals as block entropies]
    \label{thm:mpdo_tripartite_block_entropy_identifications}
    \lean{vonNeumannEntropy_tripartiteSplit_eq_blockEntropy}
    \lean{vonNeumannEntropy_traceC_eq_blockEntropy}
    \lean{vonNeumannEntropy_traceAC_eq_blockEntropy}
    \lean{vonNeumannEntropy_traceA_eq_blockEntropy}
    \leanok
    \uses{def:block_entropy}
    Let the $N$-site operator generated by $M$ be positive semidefinite, and let
    three consecutive regions have lengths $a$, $b$, and $c$, with
    $a+b+c\leq N$. After reindexing the reduced state of the first $a+b+c$
    sites as a tripartite state on $A\otimes B\otimes C$, its four entropies in
    strong subadditivity are
    \begin{align}
        S(ABC)&=S_{a+b+c}, &
        S(AB)&=S_{a+b}, \notag\\
        S(B)&=S_b, &
        S(BC)&=S_{b+c}.
        \notag
    \end{align}
\end{theorem}
\begin{proof}\leanok
    Write $\rho^{(k)}$ for the first-$k$-site reduced state, with $\cong$
    denoting equality up to a permutation of basis indices. Reindexing a
    matrix does not change its von Neumann entropy, giving the first equality
    $S(ABC)=S_{a+b+c}$ immediately. Tracing out $C$ leaves the first $a+b$
    sites, $\tr_C(\rho_{ABC})\cong\rho^{(a+b)}$, so
    $S(AB)=S_{a+b}$. Tracing out $A$ and $C$ leaves the middle $b$ sites,
    while tracing out $A$ leaves the final $b+c$ sites; translation invariance
    gives
    \begin{align}
        \tr_{A,C}(\rho_{ABC})&\cong\rho^{(b)}, \notag\\
        \tr_A(\rho_{ABC})&\cong\rho^{(b+c)}.
        \notag
    \end{align}
    Hence $S(B)=S_b$ and $S(BC)=S_{b+c}$.
\end{proof}

\begin{proposition}[Monotonicity of the mutual information]
    \label{prop:mutual_info_monotone}
    \lean{mutualInfoChain_monotone}
    \leanok
    \uses{def:mutual_info_chain}
    Let $M$ generate an MPDO whose finite-chain operator $\rho^{(N)}(M)$ is
    positive semidefinite with nonzero trace. For every block length $L$ with
    $2L+1\le N$, one has $I_L\le I_{L+1}$. The implication
    $1\le L<\lfloor N/2\rfloor\Longrightarrow 2L+1\le N$
    shows that this contains the range of
    \cite[Proposition~4.5, line~801]{Cirac2016MPDO_arXiv}.
\end{proposition}
\begin{proof}\leanok
    \uses{thm:entropy_strong_subadditivity, def:block_entropy,
        thm:mpdo_tripartite_block_entropy_identifications}
    Apply strong subadditivity to the reduced state of the first $N-L$ spins,
    split into three contiguous segments of lengths $1$, $L$, and $N-2L-1$.
    Tracing the appropriate segments yields the four block entropies
    $S_{N-L}$, $S_L$, $S_{L+1}$, and $S_{N-L-1}$, where the cyclic invariance of
    $\rho^{(N)}(M)$ identifies the entropy of a contiguous block with that of any
    block of the same length. Strong subadditivity then gives
    $S_{N-L}+S_L\le S_{L+1}+S_{N-L-1}$, which rearranges to
    $I_L\le I_{L+1}$ after subtracting $S_N$.
\end{proof}

\begin{definition}[Parity-sensitive classical MPDO]
    \label{def:parity_mpdo_counterexample}
    \lean{MPOTensor.ThermodynamicLimitCounterexample.parityTensor}
    \leanok
    \uses{def:mpo_tensor}
    Let $M$ have physical dimension two and bond dimension three, with
    \begin{align}
        M^{00}&=\diag(1,0,0), &
        M^{11}&=\diag(0,1,-1), \notag\\
        M^{01}&=0, &
        M^{10}&=0.
        \notag
    \end{align}
\end{definition}

\begin{theorem}[A periodic MPDO without a thermodynamic limit]
    \label{thm:parity_mpdo_no_thermodynamic_limit}
    \lean{MPOTensor.ThermodynamicLimitCounterexample.mpo_parityTensor_eq_diagonal,
      MPOTensor.ThermodynamicLimitCounterexample.parityTensor_isMPDO,
      MPOTensor.ThermodynamicLimitCounterexample.trace_mpo_parityTensor,
      MPOTensor.ThermodynamicLimitCounterexample.normalizedMPO_zeroConfig_of_odd,
      MPOTensor.ThermodynamicLimitCounterexample.normalizedMPO_zeroConfig_of_even,
      MPOTensor.ThermodynamicLimitCounterexample.reducedBlockState_one_zero_zero_of_odd,
      MPOTensor.ThermodynamicLimitCounterexample.reducedBlockState_one_zero_zero_of_even,
      MPOTensor.ThermodynamicLimitCounterexample.not_exists_tendsto_reducedBlockState_one_zero_zero}
    \leanok
    \uses{def:mpdo, def:parity_mpdo_counterexample}
    The tensor in Definition~\ref{def:parity_mpdo_counterexample} generates the
    positive operators
    \begin{align}
        \rho^{(N)}
        &=\ketbra{0^N}{0^N}
          +(1+(-1)^N)\ketbra{1^N}{1^N}.
        \notag
    \end{align}
    Their traces are $2+(-1)^N$. The all-zero diagonal entry of the normalized
    one-site reduced state is therefore one for odd $N$ and $1/3$ for positive
    even $N$. In particular, the normalized periodic one-site reduced states
    do not converge as $N\to\infty$.

    Thus normalized periodic marginals need not converge under the hypotheses
    of \cite[Proposition~4.5, lines~801--806]{Cirac2016MPDO_arXiv}. The mutual
    information of this same family is computed in
    \path{docs/paper-gaps/cpgsv17_mpdo_mutual_information_bound.tex}.
\end{theorem}

\begin{proof}\leanok
    A contraction around the virtual loop vanishes unless the bra and ket
    configurations agree and are both constant. The all-zero configuration
    receives weight one. The all-one configuration receives the sum of the
    two virtual-sector weights, namely $1+(-1)^N$. This gives the displayed
    diagonal form. Its coefficients are non-negative: the second coefficient
    is zero for odd $N$ and two for even $N$. Hence every positive-length
    operator is positive semidefinite, and summing its two possible diagonal
    entries gives the trace $2+(-1)^N$.

    After normalization, the all-zero weight is consequently one at odd
    lengths and $1/3$ at positive even lengths. Tracing out all but the first
    site does not change this entry, since among the configurations beginning
    with zero only the all-zero configuration has nonzero weight. The two
    subsequences of one-site marginals therefore have distinct all-zero
    entries, so the sequence cannot converge.
\end{proof}

\begin{theorem}[Mutual information without a thermodynamic limit]
    \label{thm:parity_mpdo_no_mutual_info_limit}
    \lean{MPOTensor.ThermodynamicLimitCounterexample.normalizedMPO_parityTensor_of_odd,
      MPOTensor.ThermodynamicLimitCounterexample.normalizedMPO_parityTensor_of_even,
      MPOTensor.ThermodynamicLimitCounterexample.reducedBlockState_parityTensor_of_odd,
      MPOTensor.ThermodynamicLimitCounterexample.reducedBlockState_parityTensor_of_even,
      MPOTensor.ThermodynamicLimitCounterexample.blockEntropy_parityTensor_of_odd,
      MPOTensor.ThermodynamicLimitCounterexample.blockEntropy_parityTensor_of_even,
      MPOTensor.ThermodynamicLimitCounterexample.mutualInfoChain_parityTensor_of_odd,
      MPOTensor.ThermodynamicLimitCounterexample.mutualInfoChain_parityTensor_of_even,
      MPOTensor.ThermodynamicLimitCounterexample.binEntropy_one_div_three_pos,
      MPOTensor.ThermodynamicLimitCounterexample.parityMutualInfoAfterCut,
      MPOTensor.ThermodynamicLimitCounterexample.not_exists_tendsto_parityMutualInfoAfterCut}
    \leanok
    \uses{def:mutual_info_chain, def:parity_mpdo_counterexample}
    Fix $L\geq1$. For every $N>L$, the tensor in
    Definition~\ref{def:parity_mpdo_counterexample} satisfies
    \begin{align}
        I_L^{(N)}
        &=
        \begin{cases}
            0, & N\text{ odd},\\
            h_2(1/3), & N\text{ even},
        \end{cases}
        \notag\\
        h_2(1/3)
        &=-\frac{1}{3}\log\frac{1}{3}
          -\frac{2}{3}\log\frac{2}{3}>0.
        \notag
    \end{align}
    Consequently the sequence $K\mapsto I_L^{(L+K+1)}$ does not converge,
    giving the counterexample described in
    \cite[Proposition~4.5, lines~801--806]{Cirac2016MPDO_arXiv}.
\end{theorem}

\begin{proof}\leanok
    At odd $N$, the normalized state and both of its nonempty marginals are
    pure, so all three terms in $I_L^{(N)}=S_L+S_{N-L}-S_N$ vanish. At even
    $N$, the full state and both marginals have the same two nonzero
    eigenvalues $1/3$ and $2/3$. Each of their entropies is therefore
    $h_2(1/3)$, and hence
    $I_L^{(N)}=h_2(1/3)+h_2(1/3)-h_2(1/3)=h_2(1/3)$.
    The binary entropy is strictly positive at $1/3$. Thus the odd and even
    subsequences have distinct constant values.
\end{proof}

\begin{definition}[Map on the second tensor factor]
    \label{def:mpdo_id_tensor_map}
    \lean{Matrix.idTensorMap}
    \leanok
    \uses{def:tensor_map_id}
    For a linear map $\Psi_B$ and a bipartite matrix $\rho_{AB}$, define
    \begin{align}
        (\id_A\otimes\Psi_B)(\rho_{AB})
        &=\Sigma_{B'A}\bigl((\Psi_B\otimes\id_A)
          (\Sigma_{AB}(\rho_{AB}))\bigr),
        \notag
    \end{align}
    where $\Sigma$ denotes the canonical exchange of the two tensor factors.
\end{definition}

\begin{definition}[Independent maps on both tensor factors]
    \label{def:mpdo_tensor_map_both}
    \lean{Matrix.tensorMapBoth}
    \leanok
    \uses{def:tensor_map_id, def:mpdo_id_tensor_map}
    For linear maps $\Phi_A$ and $\Psi_B$, set
    \begin{align}
        (\Phi_A\otimes\Psi_B)(\rho)
        &=(\id_{A'}\otimes\Psi_B)
          ((\Phi_A\otimes\id_B)(\rho)).
        \notag
    \end{align}
\end{definition}

\begin{lemma}[Positivity under a map on the first tensor factor]
    \label{lem:mpdo_tensor_map_id_pos}
    \lean{IsKrausCPTP.tensorMapId_posSemidef}
    \leanok
    \uses{def:is_kraus_cptp, def:tensor_map_id}
    If $\Phi_A$ is trace-preserving completely positive and
    $\rho_{AB}\geq0$, then
    $(\Phi_A\otimes\id_B)(\rho)\geq0$.
\end{lemma}

\begin{lemma}[Positivity under a map on the second tensor factor]
    \label{lem:mpdo_id_tensor_map_pos}
    \lean{IsKrausCPTP.idTensorMap_posSemidef}
    \leanok
    \uses{def:is_kraus_cptp, def:mpdo_id_tensor_map,
      lem:mpdo_tensor_map_id_pos}
    If $\Psi_B$ is trace-preserving completely positive and
    $\rho_{AB}\geq0$, then
    $(\id_A\otimes\Psi_B)(\rho)\geq0$.
\end{lemma}

\begin{lemma}[Positivity under independent maps]
    \label{lem:mpdo_tensor_map_both_pos}
    \lean{IsKrausCPTP.tensorMapBoth_posSemidef}
    \leanok
    \uses{def:mpdo_tensor_map_both, lem:mpdo_tensor_map_id_pos,
      lem:mpdo_id_tensor_map_pos}
    If $\Phi_A$ and $\Psi_B$ are trace-preserving completely positive and
    $\rho_{AB}\geq0$, then $(\Phi_A\otimes\Psi_B)(\rho)\geq0$.
\end{lemma}

\begin{theorem}[Data processing on the first tensor factor]
    \label{thm:mpdo_mutual_information_data_processing_left}
    \lean{IsKrausCPTP.mutualInformation_tensorMapId_le}
    \leanok
    \uses{def:is_kraus_cptp, def:mutual_information, def:tensor_map_id}
    Let $\rho_{AB}$ be a bipartite density operator and let $\Phi_A$ be a
    trace-preserving completely positive map whose input and output matrix
    algebras may have different dimensions. Then
    $I(A':B)_{(\Phi_A\otimes\id_B)(\rho)}\leq I(A:B)_\rho$.
\end{theorem}

\begin{proof}
    \leanok
    \uses{lem:mpdo_tensor_map_id_pos,
      thm:mutual_information_monotone_tripartite}
    Choose a rectangular Stinespring isometry for $\Phi_A$. Conjugating by
    $W=V_A\otimes\id_B$ gives
    $\omega_{A'EB}=W\rho_{AB}W^\dagger$. Since
    $W=V_A\otimes\id_B$ and $V_A^\dagger V_A=\id_A$, its marginals satisfy
    \begin{align}
        \omega_{A'E}&=V_A\rho_A V_A^\dagger, \notag\\
        \omega_B&=\rho_B.
        \notag
    \end{align}
    Entropy is preserved by the isometry $W$, and also by $V_A$ on the first
    marginal. Therefore
    \begin{align}
        S(\omega_{A'EB})&=S(\rho_{AB}), \notag\\
        S(\omega_{A'E})&=S(\rho_A), \notag\\
        S(\omega_B)&=S(\rho_B).
        \notag
    \end{align}
    Hence $I(A'E:B)_\omega=I(A:B)_\rho$. The defining Stinespring identity is
    $\tr_E(\omega_{A'EB})=(\Phi_A\otimes\id_B)(\rho_{AB})$.
    Strong subadditivity in the form $I(R:A')\leq I(R:A'E)$ proves the stated
    inequality with $R=B$.
\end{proof}

\begin{theorem}[Data processing on the second tensor factor]
    \label{thm:mpdo_mutual_information_data_processing_right}
    \lean{IsKrausCPTP.mutualInformation_idTensorMap_le}
    \leanok
    \uses{def:mpdo_id_tensor_map}
    Let $\rho_{AB}$ be a bipartite density operator and let $\Psi_B$ be a
    trace-preserving completely positive map whose input and output matrix
    algebras may have different dimensions. Then
    $I(A:B')_{(\id_A\otimes\Psi_B)(\rho)}\leq I(A:B)_\rho$.
\end{theorem}

\begin{proof}\leanok
    \uses{lem:mpdo_id_tensor_map_pos,
      thm:mpdo_mutual_information_data_processing_left}
    Exchange the two tensor factors and apply
    Theorem~\ref{thm:mpdo_mutual_information_data_processing_left}.
\end{proof}

\begin{theorem}[Data processing for bipartite mutual information]
    \label{thm:mpdo_mutual_information_data_processing}
    \lean{IsKrausCPTP.mutualInformation_tensorMapBoth_le}
    \leanok
    \uses{def:mpdo_tensor_map_both}
    Let $\rho_{AB}$ be a bipartite density operator and let $\Phi_A$ and
    $\Psi_B$ be trace-preserving completely positive maps whose input and
    output matrix algebras may have different dimensions. Then
    $I(A':B')_{(\Phi_A\otimes\Psi_B)(\rho)}\leq I(A:B)_\rho$.
\end{theorem}

\begin{proof}\leanok
    \uses{lem:mpdo_tensor_map_both_pos,
      thm:mpdo_mutual_information_data_processing_left,
      thm:mpdo_mutual_information_data_processing_right}
    Apply the two one-sided inequalities successively.
\end{proof}

\begin{definition}[Normalized state across a consecutive bipartition]
    \label{def:mpdo_split_normalized_mpo}
    \lean{MPOTensor.chainBipartitionEquiv,
      MPOTensor.bipartitionedNormalizedMPO,
      MPOTensor.bipartitionedNormalizedMPO_posSemidef,
      MPOTensor.bipartitionedNormalizedMPO_trace,
      MPOTensor.refineFirstBlock,
      MPOTensor.coarsenFirstBlock,
      MPOTensor.refineFirstBlock_isKrausCPTP,
      MPOTensor.coarsenFirstBlock_isKrausCPTP}
    \leanok
    For $N=L+K$, split a physical configuration into consecutive intervals of
    lengths $L$ and $K$, flatten each interval to a finite index, and denote the
    resulting bipartite form of the normalized periodic MPO by
    $\sigma^{(N)}_{L:K}$. The flattened local maps on the two factors are
    denoted by $\widehat{\mathcal T}_L$ and
    $\widehat{\mathcal S}_K$.
\end{definition}

\begin{lemma}[Two virtual boundaries bound the operator-Schmidt rank]
    \label{lem:mpdo_periodic_cut_operator_schmidt_rank}
    \lean{MPOTensor.periodicCutLeftFactor,
      MPOTensor.periodicCutRightFactor,
      MPOTensor.bipartitionedNormalizedMPO_hasOperatorSchmidtDecomposition,
      MPOTensor.operatorSchmidtRank_bipartitionedNormalizedMPO_le}
    \leanok
    \uses{def:mpdo_split_normalized_mpo}
    Let $M$ be an MPO tensor of operator bond dimension $D$, and let
    $N=L+K$. Opening the two virtual bonds at the consecutive cut gives
    \begin{align}
        \sigma^{(N)}_{L:K}
        &=\sum_{a,b=1}^{D} X_{ab}\otimes Y_{ab},
        \notag
    \end{align}
    and
    \begin{align}
        \operatorname{OSR}(\sigma^{(N)}_{L:K})&\leq D^2.
        \notag
    \end{align}
    For left-block words $x,x'$ and right-block words $y,y'$, the factors are
    \begin{align}
        (X_{ab})_{x,x'}
        &=\tr[\rho^{(N)}(M)]^{-1}
          \bigl(M^{x_1x'_1}\cdots M^{x_Lx'_L}\bigr)_{ab},
        \notag
    \end{align}
    and
    \begin{align}
        (Y_{ab})_{y,y'}
        &=\bigl(M^{y_1y'_1}\cdots M^{y_Ky'_K}\bigr)_{ba}.
        \notag
    \end{align}
    This is the identity
    $\tr(AB)=\sum_{a,b}A_{ab}B_{ba}$, with the normalization scalar
    absorbed into the left factor. The algebraic decomposition does not use
    positivity. With the totalized inverse convention it also remains valid
    when the trace vanishes; a nonzero trace is needed only to interpret
    $\sigma^{(N)}_{L:K}$ as the normalized physical state.

    This lemma is the two-virtual-boundary algebraic step toward the finite
    mutual-information bound asserted in CPSV16 Proposition~\texttt{PropILILp1}.
    It does not prove $I_L\leq4\log D$: the required quantum inequality
    $I(A:B)\leq2\log\operatorname{OSR}(\rho)$ remains open.
    % The analytic obstruction is tracked in issue #4295.
\end{lemma}

\begin{proof}\leanok
    Split each closed virtual word into its left and right products. Expanding
    the trace over the two exposed virtual indices gives the displayed sum of
    $D^2$ product matrices. The minimality property of the operator-Schmidt
    rank then bounds it by the number of displayed terms.
\end{proof}

\begin{lemma}[Chain mutual information across a consecutive cut]
    \label{lem:mpdo_chain_mutual_information_bipartition}
    \lean{MPOTensor.reducedBlockState_full,
      MPOTensor.bipartitionedNormalizedMPO_traceRight,
      MPOTensor.bipartitionedNormalizedMPO_traceLeft,
      MPOTensor.mutualInfoChain_eq_mutualInformation}
    \leanok
    \uses{def:mpdo_split_normalized_mpo}
    Across the consecutive cut after the first $L$ sites, the two marginals
    of $\sigma^{(N)}_{L:N-L}$ are the reduced states on $L$ and $N-L$ sites,
    up to the standard finite-index identifications. Consequently,
    $I_L(M)$ is the ordinary bipartite mutual information of
    $\sigma^{(N)}_{L:N-L}$.
\end{lemma}

\begin{proof}\leanok
    The first marginal is the defining partial trace over the last $N-L$
    sites. Periodic translation invariance identifies the second marginal
    with the reduced state on the first $N-L$ sites. Entropy is invariant
    under the finite-index identifications, and retaining the full chain
    leaves the normalized state unchanged.
\end{proof}

\begin{lemma}[Renormalization maps transfer one site across a cut]
    \label{lem:mpdo_rfp_transfer_across_cut}
    \lean{MPOTensor.tensorMapBoth_bipartitionedNormalizedMPO,
      MPOTensor.tensorMapBoth_bipartitionedNormalizedMPO_reverse}
    \leanok
    \uses{def:mpdo_split_normalized_mpo,
      thm:mpdo_global_renormalization_maps}
    Suppose that $N=(a+1)+(b+2)=(a+2)+(b+1)$. The local closure equations
    imply
    \begin{align}
        (\widehat{\mathcal T}_a\otimes\widehat{\mathcal S}_b)
        (\sigma^{(N)}_{a+1:b+2})
        &=\sigma^{(N)}_{a+2:b+1},
        \notag\\
        (\widehat{\mathcal S}_a\otimes\widehat{\mathcal T}_b)
        (\sigma^{(N)}_{a+2:b+1})
        &=\sigma^{(N)}_{a+1:b+2}.
        \notag
    \end{align}
\end{lemma}

\begin{proof}\leanok
    Fixing the word on the second interval leaves a physical closure on the
    first interval, whose virtual boundary is the fixed word evaluation.
    Apply the appropriate global closure identity to this slice. Fixing the
    resulting word on the first interval gives the closure on the second
    interval. Cyclicity of the virtual trace identifies the two cuts of the
    periodic contraction. The reverse identity is the same argument with
    $\mathcal S$ and $\mathcal T$ interchanged.
\end{proof}

\begin{theorem}[Conditional RFP mutual-information equality]
    \label{thm:mpdo_rfp_split_mutual_information_eq}
    \lean{MPOTensor.mutualInformation_bipartition_eq_of_local_transfers,
      MPOTensor.mutualInformation_bipartition_eq_of_isRFPViaTS}
    \leanok
    \uses{lem:mpdo_rfp_transfer_across_cut,
      thm:mpdo_mutual_information_data_processing}
    Let $M$ be a renormalization fixed point. If its $N$-site operator is
    positive semidefinite and has nonzero trace, then
    $I(\sigma^{(N)}_{a+1:b+2})
    =I(\sigma^{(N)}_{a+2:b+1})$.
    This is the channel calculation in
    \cite[Appendix~C, lines~1333--1341]{Cirac2016MPDO_arXiv}. For the forward
    implication of the source proposition, the required nonzero trace follows
    from horizontal canonical form, the MPDO condition, and the renormalization
    maps; simplicity is not used.
\end{theorem}

\begin{proof}\leanok
    Apply mutual-information data processing to the forward pair of local
    channels and to the reverse pair. The two transfer identities turn the
    resulting inequalities into opposite inequalities between the displayed
    mutual informations.
\end{proof}

\begin{lemma}[Nonvanishing of positive-length RFP rings]
    \label{lem:mpdo_rfp_positive_length_trace_nonzero}
    \lean{MPOTensor.firstSiteMatrix_one,
      MPOTensor.trace_mpo_ne_zero_of_isHorizontalCF_isMPDO_isRFPViaTS}
    \leanok
    \uses{def:is_rfp_via_ts}
    Let $M$ be horizontally canonical, generate positive semidefinite ring
    operators, and satisfy the local renormalization fixed-point equations.
    Then $\tr\rho^{(N)}(M)\neq0$ for every $N\geq1$.
\end{lemma}

\begin{proof}\leanok
    Horizontal canonical form gives a nonzero sector compression for one
    positive length. Positivity makes its trace nonzero. The fixed-point
    equations make the physical-trace transfer idempotent, so the trace of its
    positive powers is independent of the positive exponent.
\end{proof}

\begin{theorem}[Renormalization fixed points saturate the area law]
    \label{thm:mpdo_rfp_implies_sal}
    \lean{MPOTensor.isSAL_of_isRFPViaTS}
    \leanok
    \uses{lem:mpdo_chain_mutual_information_bipartition,
      lem:mpdo_rfp_transfer_across_cut,
      thm:mpdo_rfp_split_mutual_information_eq,
      lem:mpdo_rfp_positive_length_trace_nonzero}
    Let $M$ be a horizontally canonical matrix product density operator which
    satisfies the local renormalization fixed-point equations of
    Definition~4.1. Then $M$ saturates the area law:
    $I_L(M)=I_{L+1}(M)$ whenever
    $1\leq L<\lfloor N/2\rfloor$.
    This is the forward implication of Proposition~\texttt{propsimple} in
    Appendix~C of \cite{Cirac2016MPDO_arXiv}. Simplicity is among the ambient
    hypotheses of Theorem~4.9 but is not used in this implication.
\end{theorem}

\begin{proof}\leanok
    Write the cut as $(a+1):(b+2)$. The local channels transfer one site to
    obtain $(a+2):(b+1)$, and data processing in both directions gives equality
    of the two bipartite mutual informations. The consecutive-cut
    identification turns this into $I_L=I_{L+1}$. The preceding lemma supplies
    the normalization at every positive chain length.
\end{proof}

\begin{theorem}[Renormalization fixed points have ZCL and SAL]
    \label{thm:mpdo_rfp_implies_zcl_and_sal}
    \lean{MPOTensor.isSourceZCL_and_isSAL_of_isRFPViaTS}
    \leanok
    \uses{thm:rfp_via_ts_source_zcl,
      lem:mpdo_rfp_positive_length_trace_nonzero,
      thm:mpdo_rfp_implies_sal}
    Let $M$ be a horizontally canonical matrix product density operator which
    satisfies the local renormalization fixed-point equations of
    Definition~4.1. Then $M$ has source zero correlation length and saturates
    the area law. This is Proposition~\texttt{propsimple} in Appendix~C of
    \cite{Cirac2016MPDO_arXiv}, lines~1333--1341.
\end{theorem}

\begin{proof}\leanok
    The preceding nonvanishing lemma at length one rules out a zero
    physical-trace transfer. Trace preservation of the refinement map then
    gives source zero correlation length, while
    Theorem~\ref{thm:mpdo_rfp_implies_sal} gives saturation of the area law.
\end{proof}

\begin{lemma}[Nonnegativity of block entropies]
    \label{lem:block_entropy_nonneg}
    \lean{blockEntropy_nonneg}
    \leanok
    \uses{def:block_entropy, thm:entropy_nonneg}
    For a normalized positive semidefinite finite-chain operator, every block entropy
    is non-negative: $0\le S_m$.
\end{lemma}
\begin{proof}\leanok
    The reduced state of a block is a density matrix, and the von Neumann entropy of
    a density matrix is non-negative.
\end{proof}

\begin{lemma}[Nonnegativity of the mutual information]
    \label{lem:mutual_info_chain_nonneg}
    \lean{mutualInfoChain_nonneg}
    \leanok
    \uses{def:mutual_info_chain, thm:entropy_strong_subadditivity, lem:block_entropy_nonneg}
    Let $M$ generate an MPDO whose finite-chain operator $\rho^{(N)}(M)$ is positive
    semidefinite with nonzero trace. For every block length $L\le N$, one has
    $0\le I_L$.
\end{lemma}
\begin{proof}\leanok
    \uses{thm:entropy_strong_subadditivity, def:block_entropy, lem:block_entropy_nonneg}
    Apply strong subadditivity with a trivial middle subsystem (equivalently,
    subadditivity) to the bipartition of the chain into the first $L$ and last
    $N-L$ spins: $S_N+S_0\le S_L+S_{N-L}$. Since $S_0\ge0$
    (Lemma~\ref{lem:block_entropy_nonneg}), rearranging gives
    $I_L=S_L+S_{N-L}-S_N\ge S_0\ge0$.
\end{proof}

\begin{lemma}[Symmetry of the mutual information across the cut]
    \label{lem:mutual_info_chain_symm}
    \lean{mutualInfoChain_symm}
    \leanok
    \uses{def:mutual_info_chain}
    For every block length $L\le N$, the mutual information of the first $L$ spins
    equals that of their complement: $I_L=I_{N-L}$.
\end{lemma}
\begin{proof}\leanok
    \uses{def:mutual_info_chain}
    Both sides equal $S_L+S_{N-L}-S_N$ by the symmetric definition
    $I_L=S_L+S_{N-L}-S_N$, using $N-(N-L)=L$.
\end{proof}

\begin{lemma}[Positive rank of a normalized state]
    \label{lem:trace_one_positive_rank}
    \lean{Matrix.PosSemidef.rank_pos_of_trace_one}
    \leanok
    A positive semidefinite matrix of trace one has positive rank.
\end{lemma}
\begin{proof}\leanok
    If its rank were zero, every eigenvalue would vanish, contradicting that
    their sum is the trace, which equals one.
\end{proof}

\begin{lemma}[Zero entropy in rank one]
    \label{lem:rank_one_entropy_zero}
    \lean{vonNeumannEntropy_eq_zero_of_rank_le_one}
    \leanok
    \uses{lem:trace_one_positive_rank, thm:entropy_le_log_rank,
        thm:entropy_nonneg}
    A positive semidefinite matrix of trace one and rank at most one has zero
    von Neumann entropy.
\end{lemma}
\begin{proof}\leanok
    \uses{lem:trace_one_positive_rank, thm:entropy_le_log_rank,
        thm:entropy_nonneg}
    Its rank is positive by Lemma~\ref{lem:trace_one_positive_rank}, hence it is
    exactly one. The rank bound gives $S\le\log 1=0$, while entropy is
    non-negative.
\end{proof}

\begin{lemma}[Entropy of the empty block]
    \label{lem:block_entropy_zero}
    \lean{blockEntropy_zero}
    \leanok
    \uses{def:block_entropy, lem:rank_one_entropy_zero}
    For a normalized positive semidefinite finite-chain operator, the block entropy of
    the empty block vanishes: $S_0=0$.
\end{lemma}
\begin{proof}\leanok
    \uses{lem:rank_one_entropy_zero}
    The reduced state of zero spins is one-dimensional, so its rank is at most
    one. Lemma~\ref{lem:rank_one_entropy_zero} gives $S_0=0$.
\end{proof}

\begin{lemma}[Empty-block mutual information vanishes]
    \label{lem:mutual_info_chain_zero}
    \lean{mutualInfoChain_zero}
    \leanok
    \uses{def:mutual_info_chain, lem:block_entropy_zero}
    For a normalized positive semidefinite finite-chain operator, the mutual information
    of the empty block vanishes: $I_0=0$.
\end{lemma}
\begin{proof}\leanok
    \uses{lem:block_entropy_zero}
    By definition, $I_0=S_0+S_N-S_N=S_0$, which is zero by
    Lemma~\ref{lem:block_entropy_zero}.
\end{proof}

\begin{definition}[Ancillary trace on a block]
    \label{def:block_ancillary_trace}
    \lean{MPOTensor.blockSpinAncillaEquiv,
      MPOTensor.blockSpinAncillaEquiv_symm_apply,
      MPOTensor.blockSpinAncillaEquiv_symm_decode,
      MPOTensor.blockAncillaryTraceMap}
    \leanok
    \uses{def:mpdo_ancillary_trace_map}
    For a block of $L$ spin--ancilla pairs, regroup
    \begin{align}
        (\C^d\otimes\C^{d_K})^{\otimes L}
        &\cong(\C^d)^{\otimes L}
          \otimes(\C^{d_K})^{\otimes L}
        \notag
    \end{align}
    and trace the ancillary factor. Denote this map by
    $\Tr_{\mathrm{anc},L}$. It is the blockwise analog of the
    one-site ancillary trace in
    Definition~\ref{def:mpdo_ancillary_trace_map}.
\end{definition}

\begin{theorem}[Block ancillary trace is a channel]
    \label{thm:block_ancillary_trace_channel}
    \lean{MPOTensor.blockAncillaryTraceMap_isKrausCPTP}
    \leanok
    \uses{def:block_ancillary_trace, def:is_kraus_cptp}
    The map $\Tr_{\mathrm{anc},L}$ is trace-preserving and
    completely positive.
\end{theorem}

\begin{proof}\leanok
    \uses{def:block_ancillary_trace}
    The reindexing equivalence is a channel, as is the partial trace over the
    ancillary factor. Their composition is therefore trace-preserving and
    completely positive.
\end{proof}

\begin{theorem}[Coefficients of the ancillary traces]
    \label{thm:block_ancillary_trace_coefficients}
    \lean{MPOTensor.blockAncillaryTraceMap_apply,
      MPOTensor.tensorMapBoth_blockAncillaryTraceMap_apply}
    \leanok
    \uses{def:block_ancillary_trace}
    If $X$ is an operator on a spin--ancilla block, then for
    $i,j\in[d]^L$,
    \begin{align}
        \Tr_{\mathrm{anc},L}(X)_{i,j}
        &=\sum_{\kappa\in[d_K]^L}X_{(i,\kappa),(j,\kappa)}.
        \notag
    \end{align}
    For $i,j\in[d]^L$ and $x,y\in[d]^K$, applying the ancillary traces to
    both blocks across the cut $L\mid K$ gives
    \begin{align}
        \bigl[(\Tr_{\mathrm{anc},L}\otimes
          \Tr_{\mathrm{anc},K})(X)\bigr]_{(i,x),(j,y)}
        &=\sum_{\kappa\in[d_K]^L}\sum_{\lambda\in[d_K]^K}
          X_{((i,\kappa),(x,\lambda)),((j,\kappa),(y,\lambda))}.
        \notag
    \end{align}
\end{theorem}

\begin{proof}\leanok
    \uses{def:block_ancillary_trace}
    After regrouping the spin and ancillary factors, the definition of the
    partial trace gives the first coefficient identity. Applying this identity
    to both tensor factors gives the displayed iterated sum; the two finite
    sums may be interchanged.
\end{proof}

\begin{theorem}[Local purification as a block-channel image]
    \label{thm:lpdo_block_channel_image}
    \lean{MPOTensor.blockAncillaryTraceMap_pureState_purificationTensor,
      MPOTensor.trace_purificationDensity_eq_pureState,
      MPOTensor.bipartitionedNormalizedMPO_eq_tensorMapBoth_blockAncillaryTraceMap}
    \leanok
    \uses{def:block_ancillary_trace,
      def:mpdo_purification_tensor, def:mpdo_purification_density,
      def:mpdo_split_normalized_mpo}
    Fix a local-purification representation
    \begin{align}
        M^{ij}
        &=\left(\sum_{k=0}^{d_K-1}
          A^{(i,k)}\otimes\overline{A^{(j,k)}}\right)_{e,e},
        \notag
    \end{align}
    with purifying MPS tensor $\widetilde A$ of physical dimension $d d_K$
    and bond dimension $D'$. For every $N=L+K$, tracing all ancillary
    degrees of freedom gives
    \begin{align}
        \rho^{(N)}(M)
        &=\Tr_{\mathrm{anc},N}
          \left(\ketbra{V^{(N)}(\widetilde A)}
            {V^{(N)}(\widetilde A)}\right), \notag\\
        \tr\rho^{(N)}(M)
        &=\left\lVert V^{(N)}(\widetilde A)\right\rVert^2.
        \notag
    \end{align}

    Define the normalized-form purifying operator by
    \begin{align}
        \pi_{\widetilde A}^{(N)}
        &:=\left\lVert V^{(N)}(\widetilde A)\right\rVert^{-2}
          \ketbra{V^{(N)}(\widetilde A)}
            {V^{(N)}(\widetilde A)},
        \notag
    \end{align}
    where the inverse scalar is taken with the convention $0^{-1}=0$. After
    reindexing across the cut $L\mid K$,
    \begin{align}
        \sigma^{(N)}_{L:K}(M)
        &=\left(\Tr_{\mathrm{anc},L}\otimes
          \Tr_{\mathrm{anc},K}\right)
          \left(\left(\pi_{\widetilde A}^{(N)}\right)_{L:K}\right).
        \notag
    \end{align}
    This identity also holds when the common trace vanishes, in which case both
    normalized-form operators are zero.

    This is the finite-chain channel identity from the purification preceding
    equation~(4) of \cite{Wolf2008AreaLaws}, under the explicit local
    purification of \cite[Section~4.3]{Cirac2016MPDO_arXiv}.
\end{theorem}

\begin{proof}\leanok
    \uses{thm:block_ancillary_trace_coefficients,
      thm:block_ancillary_trace_channel,
      thm:mpo_eq_purification_density}
    The coefficient formula for one block gives
    \begin{align}
        &\Tr_{\mathrm{anc},N}
          \left(\ketbra{V^{(N)}(\widetilde A)}
            {V^{(N)}(\widetilde A)}\right)_{\sigma,\tau}
        \notag\\
        &\quad
        =\sum_{\kappa\in[d_K]^N}
          V^{(N)}(\widetilde A)_{(\sigma,\kappa)}
          \overline{V^{(N)}(\widetilde A)_{(\tau,\kappa)}}
        =\rho^{(N)}(M)_{\sigma,\tau}.
        \notag
    \end{align}
    Trace preservation gives equality of the two normalizing traces. Splitting
    $\kappa$ into its left and right restrictions changes the last sum into
    the two independent sums in
    Theorem~\ref{thm:block_ancillary_trace_coefficients}, which proves the
    normalized identity across $L\mid K$.
\end{proof}

\begin{theorem}[LPDOs admit a block-channel representation]
    \label{thm:lpdo_exists_block_channel_image}
    \lean{MPOTensor.IsLPDO.%
      exists_bipartitionedNormalizedMPO_eq_blockAncillaryTrace}
    \leanok
    \uses{def:lpdo, thm:lpdo_block_channel_image}
    Let $M$ be an LPDO. For every decomposition $N=L+K$, there exist an
    ancillary dimension $d_K$, a purifying bond dimension $D'$, and a
    local-purification family $A$. Let $\widetilde A$ be its associated
    purifying MPS tensor. Then
    \begin{align}
        \sigma^{(N)}_{L:K}(M)
        &=\left(\Tr_{\mathrm{anc},L}\otimes
          \Tr_{\mathrm{anc},K}\right)
          \left(\left(\pi_{\widetilde A}^{(N)}\right)_{L:K}\right).
        \notag
    \end{align}
\end{theorem}

\begin{proof}\leanok
    \uses{def:lpdo, thm:lpdo_block_channel_image}
    Choose a local-purification representation of $M$ and apply
    Theorem~\ref{thm:lpdo_block_channel_image} to its purifying tensor.
\end{proof}

\begin{theorem}[LPDOs admit a cut-uniform block-channel representation]
    \label{thm:lpdo_exists_cut_uniform_block_channel_image}
    \lean{MPOTensor.IsLPDO.%
      exists_forall_bipartitionedNormalizedMPO_eq_blockAncillaryTrace}
    \leanok
    \uses{def:lpdo, thm:lpdo_block_channel_image}
    Let $M$ be an LPDO. There exist an ancillary dimension $d_K$, a
    purifying bond dimension $D'$, and a local-purification family $A$, all
    independent of the cut, such that for every decomposition $N=L+K$,
    \begin{align}
        \sigma^{(N)}_{L:K}(M)
        &=\left(\Tr_{\mathrm{anc},L}\otimes
          \Tr_{\mathrm{anc},K}\right)
          \left(\left(\pi_{\widetilde A}^{(N)}\right)_{L:K}\right),
        \notag
    \end{align}
    where $\widetilde A$ is the purifying MPS tensor associated with $A$.
\end{theorem}

\begin{proof}\leanok
    \uses{def:lpdo, thm:lpdo_block_channel_image}
    Choose one local-purification representation of $M$. For an arbitrary
    decomposition $N=L+K$, apply
    Theorem~\ref{thm:lpdo_block_channel_image} to the same purifying tensor.
\end{proof}

\begin{theorem}[Mutual-information bound for a local-channel image]
    \label{thm:local_channel_image_mutual_info_bound}
    \lean{MPOTensor.mutualInfoChain_le_of_bipartitioned_channel_image}
    \leanok
    \uses{def:mutual_info_chain, def:mpdo_split_normalized_mpo,
      def:mpdo_tensor_map_both, def:is_kraus_cptp, def:doubled_tensor}
    Let $M$ be an MPO tensor whose generated $N$-site operator is positive
    semidefinite. Let $A$ be an MPS tensor of bond dimension $D'\geq1$ whose
    $N$-site vector is nonzero. Suppose that, across the cut $L\mid(N-L)$, the
    normalized state generated by $M$ is obtained from the normalized pure state
    of $A$ by trace-preserving completely positive maps $\Phi_L$ and $\Phi_R$ on
    the two blocks:
    \begin{align}
        \sigma^{(N)}_{L:N-L}(M)
        &=(\Phi_L\otimes\Phi_R)
          \left(\frac{\ketbra{V^{(N)}(A)}{V^{(N)}(A)}}
              {\lVert V^{(N)}(A)\rVert^2}\right)_{L:N-L}.
        \notag
    \end{align}
    Then $I_L(M)\leq4\log D'$.
\end{theorem}
\begin{proof}\leanok
    \uses{thm:mpdo_mutual_information_data_processing,
      prop:pure_mutual_info_bound}
    Write $\sigma_A^{(N)}$ for the normalized pure state of $A$ in the
    $L\mid(N-L)$ bipartition. Then
    \begin{align}
        I_L(M)
        &=I\!\left((\Phi_L\otimes\Phi_R)(\sigma_A^{(N)})\right)
        \leq I(\sigma_A^{(N)})
        \leq4\log D'.
        \notag
    \end{align}
    The equality is the channel-image hypothesis, the first inequality is data
    processing, and the last inequality is
    Proposition~\ref{prop:pure_mutual_info_bound}.
\end{proof}

\begin{theorem}[Mutual-information bound for a locally purified tensor]
    \label{thm:lpdo_mutual_info_bound}
    \lean{MPOTensor.purifyingBondDimension_pos_of_trace_ne_zero,
      MPOTensor.mutualInfoChain_le_of_lpdoWitness,
      MPOTensor.IsLPDO.exists_mutualInfoChain_le_purifyingBond}
    \leanok
    \uses{def:lpdo, def:mutual_info_chain,
      thm:lpdo_block_channel_image}
    Fix a local-purification representation with purifying MPS bond dimension
    $D'$. If the $N$-site operator is positive semidefinite and its purifying
    pure state is nonzero, then $I_L(M)\leq4\log D'$ for every $L\leq N$.

    Consequently, if $M$ is an LPDO, $N\geq1$, and
    $\tr\rho^{(N)}(M)\neq0$, then its local-purification representation
    supplies a bond dimension $D'$ satisfying this estimate. The quantity
    $D'$ is the bond dimension of the purifying MPS, not the MPO bond
    dimension $D$.
\end{theorem}

\begin{proof}\leanok
    \uses{thm:lpdo_block_channel_image,
      thm:block_ancillary_trace_channel,
      thm:local_channel_image_mutual_info_bound,
      thm:lpdo_is_mpdo,
      thm:mpo_eq_purification_density}
    By Theorem~\ref{thm:lpdo_block_channel_image},
    \begin{align}
        I_L(M)
        &=I\!\left((\Tr_{\mathrm{anc},L}\otimes
            \Tr_{\mathrm{anc},N-L})
          \left(\left(\pi_{\widetilde A}^{(N)}\right)_{L:N-L}\right)\right)
        \notag\\
        &\leq I\!\left(\left(\pi_{\widetilde A}^{(N)}\right)_{L:N-L}\right)
        \leq4\log D'.
        \notag
    \end{align}
    The first inequality is data processing and the second is the pure-state
    boundary estimate. Nonvanishing of the purifying pure state implies
    $D'>0$. For an LPDO, positivity follows from
    Theorem~\ref{thm:lpdo_is_mpdo}, while the trace identity in
    Theorem~\ref{thm:lpdo_block_channel_image} transfers the nonzero-trace
    hypothesis to the purifying pure state.
\end{proof}
